% theme: applications_with_lines
\MajorQuestionLesson{直線との利用問題}
\SetRowSpace{4mm}

\TypeQuestion{図のように、放物線 $y=ax^2$ と直線 $\ell:y=mx+6$ が2点 $\mathrm{A},\mathrm{B}$ で交わっている。点 $\mathrm{A}$ の座標は $(-2,4)$ で、点 $\mathrm{B}$ の $x$ 座標は正である。また、点 $\mathrm{O}$ は原点である。次の問いに答えなさい。}{line_application_coefficients}
\QQ{定数 $a$ の値を求めなさい。}{}
\QQ{定数 $m$ の値を求めなさい。}{}
\QQ{点 $\mathrm{B}$ の座標を求めなさい。}{}
\QQ{$\Tri{OAB}$ の面積を求めなさい。}{}

\TypeQuestion{図のように、放物線 $y=\dfrac{1}{2}x^2$ と直線 $\ell$ が2点 $\mathrm{A},\mathrm{B}$ で交わり、点 $\mathrm{A},\mathrm{B}$ の $x$ 座標はそれぞれ $-2,4$ である。直線 $m$ は直線 $\ell$ と平行で、放物線と2点 $\mathrm{C},\mathrm{D}$ で交わっている。点 $\mathrm{C}$ の座標を $(-4,p)$ とし、点 $\mathrm{D}$ の $x$ 座標は正である。次の問いに答えなさい。}{line_application_parallel}
\QQ{点 $\mathrm{A},\mathrm{B}$ の座標をそれぞれ求めなさい。}{}
\QQ{直線 $\ell$ の式を求めなさい。}{}
\QQ{$p$ の値と、直線 $m$ の式を求めなさい。}{}
\QQ{点 $\mathrm{D}$ の座標を求めなさい。}{}

\LessonPageBreak

\TypeQuestion{図のように、放物線 $y=x^2$ 上に点 $\mathrm{A}(-2,4)$ と、$x$ 座標が正である点 $\mathrm{B}$ がある。点 $\mathrm{B}$ の $x$ 座標を $t$ とし、2点 $\mathrm{A},\mathrm{B}$ を通る直線を $\ell$ とする。また、直線 $\ell$ と $y$ 軸との交点を $\mathrm{C}$、原点を $\mathrm{O}$ とする。次の問いに答えなさい。}{line_application_parameter}
\QQ{点 $\mathrm{B}$ の座標を $t$ を用いて表しなさい。}{}
\QQ{直線 $\ell$ の式を $t$ を用いて表しなさい。}{}
\QQ{点 $\mathrm{C}$ の座標を $t$ を用いて表しなさい。}{}
\QQ{$\Tri{OAB}$ の面積が $24$ であるとき、$t$ の値と点 $\mathrm{B}$ の座標を求めなさい。}{}

\TypeQuestion{図のように、放物線 $y=\dfrac{1}{2}x^2$ 上に3点 $\mathrm{A}(-4,8),\mathrm{B}(p,2),\mathrm{C}(q,8)$ があり、$p>0,\ q>0$ である。2点 $\mathrm{A},\mathrm{B}$ を通る直線を $\ell$、2点 $\mathrm{B},\mathrm{C}$ を通る直線を $m$ とする。また、直線 $\ell,m$ と $y$ 軸との交点をそれぞれ $\mathrm{D},\mathrm{E}$ とする。次の問いに答えなさい。}{line_application_two_lines}
\QQ{$p,q$ の値を求めなさい。}{}
\QQ{直線 $\ell,m$ の式をそれぞれ求めなさい。}{}
\QQ{点 $\mathrm{D},\mathrm{E}$ の座標をそれぞれ求めなさい。}{}
\QQ{$\Tri{BDE}$ の面積を求めなさい。}{}
